\documentclass[11pt,a4paper]{amsart}
\usepackage[T1]{fontenc}
\usepackage{amssymb}
\usepackage[hmargin=1in,vmargin=.9in]{geometry}
\usepackage[hidelinks]{hyperref}
\hypersetup{pdftitle={A four-vertex counterexample for one-arc Bayesian network neighbors}}
\DeclareFontShape{OT1}{cmr}{bx}{n}{<-> cmbx10}{}
\DeclareMathSizes{8}{8}{8}{8}
\DeclareMathSizes{9}{9}{8}{8}
\DeclareMathSizes{10.95}{10.95}{8}{8}
\newtheorem{theorem}{Theorem}
\title{A four-vertex counterexample for one-arc Bayesian network neighbors}
\date{}
\begin{document}
\begin{abstract}
A nearest graph in d-separation query distance need not be one
arc edit from a supplied directed acyclic graph. A four-vertex example has global
distance one and best positive one-edit distance two. Adding isolated
vertices extends the gap to every larger size. A Markov-equivalent
reorientation of the input does have a single edit to the witness.
The counterexample therefore addresses the fixed-input interpretation.
\end{abstract}
\maketitle

\section{Distance and statement}

For a DAG \(D\) on vertex set \(V\), its signature \(I(D)\)
contains \((\{u,v\},S)\) for distinct endpoints \(u,v\) and
\(S\subseteq V\setminus\{u,v\}\) that d-separates them.
A path is active given \(S\) if no internal noncollider is in
\(S\), and each internal collider lies in \(S\) or has a
descendant in \(S\). The distance \(d(D,H)\) is the size of the
symmetric difference of the two signatures. A nearest neighbor
has minimum positive distance, so Markov-equivalent DAGs at
distance zero are excluded.

Conjecture 2 of Harviainen, Parviainen and Sharma \cite[p.~2065]{hps}
asks whether such a neighbor can be reached by a single arc
insertion, deletion or reversal that preserves acyclicity. The
opening paragraph of their Section 3.2 describes this as a
transformation of the given DAG. We take the input orientation
as supplied.

\begin{theorem}
On four vertices there is a DAG whose nearest non-equivalent
alternative has d-separation distance 1, but every one-arc
edit has distance either 0 or at least 2.
For every \(n\ge4\), one can obtain a global minimum positive
distance \(2^{n-4}\) and best positive one-edit distance \(2^{n-3}\).
\end{theorem}

\section{The four-vertex witness}

Use vertices \(0,1,2,3\) and the two DAGs
\[
 D=\{0\to1,\ 0\to2,\ 0\to3,\ 1\to2\},\qquad
 H=\{1\to0,\ 2\to0,\ 0\to3\}.
\]
Adjacent endpoints cannot be d-separated. In \(D\) the only
nonadjacent pairs are \(\{1,3\},\{2,3\}\). Every path to
vertex 3 passes through noncollider 0, and the two-step route
using only that intermediate vertex is active when 0 is not
conditioned on. Thus the only four separations of \(D\) are
\[
 1\perp_d3\mid\{0\},\quad1\perp_d3\mid\{0,2\},\quad
 2\perp_d3\mid\{0\},\quad2\perp_d3\mid\{0,1\}.
\]
The unique corresponding paths in \(H\) again have 0 as a
noncollider. The additional nonadjacent pair \(\{1,2\}\)
has the sole path \(1\to0\leftarrow2\). Its collider is activated
by conditioning on 0 or its descendant 3. Hence \(H\) adds only
the separation \(1\perp_d2\mid\varnothing\). The distance is one
and therefore is globally minimal among positive distances.

\section{Every one-arc edit}

Only the pairs \(\{1,3\},\{2,3\}\) can be added. Each
new adjacency removes both prior separations of its endpoint
pair. This forces distance at least two in either direction.

Each deletion gives the following new separations on its endpoints.
\[
\begin{array}{c|l}
\text{deleted arc}&\text{conditioning sets with a new separation}\\
0\to1&\varnothing,\ \{3\}\\
0\to2&\{1\},\ \{1,3\}\\
0\to3&\varnothing,\ \{1\},\ \{2\},\ \{1,2\}\\
1\to2&\{0\},\ \{0,3\}
\end{array}
\]
For the first row the remaining route is a collider through 2.
For the second it is a chain through 1. The third isolates
vertex 3, and the fourth leaves a fork through 0. In the
fourth case the two additional separations are the only changes:
the edited graph is a fork from 0 to all three leaves and
the previous separations remain unchanged.

The reversal of \(0\to2\) is cyclic, since \(0\to1\to2\)
remains. The three other reversals are acyclic. After any of
them the skeleton and the set of unshielded colliders are
unchanged; the latter is empty in all four DAGs. The standard
equivalence criterion \cite{vp,chick95} therefore gives distance
zero in each case. This exhausts the admissible edits. It also
shows that the best positive distance among one-edit neighbors
is exactly two.

On two vertices the only distinct signatures arise with and
without an arc, so a single edit always suffices. On three
vertices, a DAG with two arcs has a path skeleton and exactly
one separation. Adding its missing arc in an acyclic orientation
reaches distance one. A complete DAG similarly reaches distance
one by a deletion. A DAG with at most one arc has both possible
separations for every nonadjacent pair. Any alternative with a
new adjacency therefore differs on at least two queries. A
one-arc alternative without a new adjacency is equivalent or
empty; deleting its arc changes two queries. Adding an arc to
the empty graph also changes exactly two. Thus four vertices
are necessary for failure.

\section{Adding isolated vertices}

Add \(n-4\) isolated vertices to both graphs. If all the new
vertices stay isolated in a comparison, each original query
is replicated \(2^{n-4}\) times by arbitrary choices of
additional conditioning vertices. The queries involving an
isolated endpoint do not change. Any alternative that creates
an arc incident to an old isolated vertex destroys the
separation of its endpoints for all \(2^{n-2}\) possible
conditioning sets. This cost is larger than the distance
achieved by the padded \(H\) and the best original edit.
The padded \(H\) is therefore still nearest at positive
distance, now \(2^{n-4}\). The least positive one-edit
distance is \(2\cdot2^{n-4}\) by the same argument.

\section{A different question about equivalence classes}

The DAG
\[
 D'=\{1\to0,\ 2\to0,\ 0\to3,\ 1\to2\}
\]
has the skeleton of \(D\) and no unshielded collider. It is
Markov equivalent to \(D\), and deleting \(1\to2\) from it
gives \(H\). Thus allowing free reorientation before the
edit changes the answer for this example.

More generally, Chickering's transformation theorem
\cite[Thm.~4]{chick02} gives an attributed implication for nested
signatures. From the graph with the larger independence model,
covered reversals and arc additions build the other graph.
The pairwise signatures suffice to test nesting of independence
models, because d-separation of two vertex sets is equivalent
to pairwise d-separation between them with the same conditions.
Along the transformation the signature is nonincreasing by
set inclusion. If the second class is nearest to the first
and there were more than one strict change, the first changed
class would be closer at positive distance. Thus exactly one
strict change is needed, with only covered reversals before
and after it. The same applies starting at the other end of
a nearest pair by reversing the sequence and using the last
strict step. In particular distance-one nearest pairs give no
obstruction to the class-level question.

The finite class-level census finds no failure through five
vertices when free reorientation is allowed. This does not
establish the general class-level claim. The source's chain
calculation has a different restriction, and the undirected
conjecture is outside the result.

\section{Verification and disclosure}

Internally reviewed, not peer reviewed. The exact path replay
checks all 24 queries on the four-vertex examples and every
possible one-arc edit. The independent Python class census
uses ancestral moralization and checks every labeled DAG
through five vertices. No proof-assistant formalization is
included.

Under human direction, ChatGPT Astra 6.2 (OpenAI) generated the
discovery proof and its research packet. OpenAI Codex independently
checked the mathematics. GPT-6.2 Sol (OpenAI) checked the literature
and prepared the manuscript, exact replay, independent class
census and illustrations. DannyExperiments maintains the repository
and is the preferred citation identity. No AI system is proposed
as an author.

\begin{thebibliography}{9}
\raggedright
\bibitem{hps}
J. Harviainen, P. Parviainen and V. S. Sharma,
\emph{Learning Bayesian and Markov Networks with an Unreliable
Oracle}, Proc. UAI, PMLR \textbf{337} (2026), 2059--2074.
\href{https://proceedings.mlr.press/v337/harviainen26a.html}{Proceedings record}.
\bibitem{vp}
T. S. Verma and J. Pearl, \emph{Equivalence and Synthesis of
Causal Models}, Proc. UAI (1990), 220--227.
\href{https://ftp.cs.ucla.edu/tech-report/1991-reports/910020.pdf}{Author report}.
\bibitem{chick95}
D. M. Chickering, \emph{A Transformational Characterization of
Equivalent Bayesian Network Structures}, Proc. UAI (1995), 87--98.
\href{https://maxchickering.com/publications/uai95.pdf}{Author PDF}.
\bibitem{chick02}
D. M. Chickering, \emph{Optimal Structure Identification With
Greedy Search}, Journal of Machine Learning Research \textbf{3}
(2002), 507--554.
\href{https://jmlr.org/papers/v3/chickering02b.html}{Journal record}.
\end{thebibliography}
\end{document}
